% answers_short.tex —— 2019 年·北京卷·文科 · 简短答案速查
% 由 tools/make_short_answers.py 自 answers.tex 的 \gksolhead/\gkanswer 提取，
% 并与 2019/2019_数学真题_北京卷文_答案要点.md 逐题比对后生成，勿手改（改源 MD 后重跑）。

\begin{gknotice}
\textbf{使用说明}：本节只给最终答案便于快速核对，完整推导见解析版。答案为编写者独立推导，\textbf{非官方标准答案}。
\end{gknotice}

\gkpart{选择题}
\begin{center}\begin{tabular}{ccccccccc}
\toprule
题号 & 1 & 2 & 3 & 4 & 5 & 6 & 7 & 8 \\
\midrule
答案 & C & D & A & B & D & C & A & B \\
\bottomrule
\end{tabular}\end{center}

\gkpart{填空题}
\begin{itemize}[leftmargin=2.6em]
  \item[\textbf{9.}] $8$
  \item[\textbf{10.}] $-3$；$1$
  \item[\textbf{11.}] $(x - 1)^2 + y^2 = 4$
  \item[\textbf{12.}] $28$
  \item[\textbf{13.}] ②③ $\implies$ ①（或 ①③ $\implies$ ②）
  \item[\textbf{14.}] ① $130$；② $15$
\end{itemize}

\gkpart{解答题}
\begin{itemize}[leftmargin=2.6em]
  \item[\textbf{15.}] (1) $b = 7, c = 5$；(2) $\dfrac{3\sqrt{3}}{14}$
  \item[\textbf{16.}] (1) $a_n = 2n - 12$；(2) $-30$
  \item[\textbf{17.}] (1) $400$ 人；(2) $\dfrac{1}{25}$（或 $0.04$）；(3) 能认为有变化（或不能认为有变化），理由见解析
  \item[\textbf{18.}] (1) 证明见解析；(2) 证明见解析；(3) 存在点 $F$（为棱 $PB$ 的中点），使得 $CF \parallel$ 平面 $PAE$，理由见解析
  \item[\textbf{19.}] (1) $\dfrac{x^2}{2} + y^2 = 1$；(2) 证明见解析，直线 $l$ 经过定点 $(0, 0)$
  \item[\textbf{20.}] (1) $y = x$ 或 $y = x - \dfrac{64}{27}$；(2) 证明见解析；(3) $a = -3$
\end{itemize}
